% ------------------------------------------------------------
% CHAPTER 21
% ------------------------------------------------------------

\chapter{Toward a General Theory of Discrete Field Sampling}

The classical Shannon sampling theorem occupies a central place in signal processing. It states that a bandlimited function can be reconstructed exactly from uniform samples taken at twice its highest frequency component. This theorem depends on linearity, translation-invariance, and ideal low-pass behavior. Delta–sigma modulation appears, at first glance, to have no direct relation to Shannon’s theory. It is nonlinear, it relies on feedback rather than interpolation, it produces quantization noise with highly structured spectra, and its sampling behavior is governed by dynamic error correction rather than linear projection. Nevertheless, delta–sigma modulation can be understood as part of a far more general theory of field sampling, one that transcends the classical Fourier picture. In this chapter, we develop such a theory by building directly on the RSVP field-theoretic framework.

In RSVP, a field is not a function defined over a static domain, but a dynamic configuration of the scalar, vector, and entropy fields. Sampling such a field corresponds not merely to evaluating it at discrete points, but to interacting with its local structure through a measurement operator that inevitably injects entropy. Delta–sigma modulation provides a canonical example of this interaction. The quantizer enforces discrete symbolic collapse at specific times, which corresponds to the evaluation of \(\Phi(t)\) under severe constraints. The feedback loop counteracts this collapse by reconstructing the field’s local curvature through repeated integration. In this way, delta–sigma modulation does not sample the field in the Shannon sense; it reconstructs it through a repeated relaxation of the RSVP dynamics.

To formalize this viewpoint, let the underlying field be a function \( \Phi : \mathbb{R} \to \mathbb{R} \), which may or may not be bandlimited. At a given sampling time \(t_n\), the quantizer produces a symbol
\[
y[n] = Q(\Phi(t_n)).
\]
This symbol corresponds to a discrete observation obtained by collapsing the continuous field to one of finitely many representational states. In classical sampling theory, this collapse would represent a linear projection. In delta–sigma modulation, however, the collapse injects entropy into the state of the system. The entropy field absorbs this injection by generating localized perturbations, which are subsequently smoothed by the vector field \(v\) corresponding to the converter’s loop filter.

The reconstruction of the field is performed not through direct interpolation, but through the evolution of the state variable \(x[n]\), which satisfies the recurrence
\[
x[n+1] = A x[n] + B u[n] - C y[n].
\]
This recurrence attempts to approximate the scalar field \(\Phi\) by integrating the discrepancy between the physical input \(u[n]\) and the quantizer output. In this sense, delta–sigma modulation implements a form of discrete variational inference: it seeks the trajectory of \(x[n]\) that minimizes the perceived discrepancy between the continuous field and its symbolic representation. The modulator’s output \(y[n]\) therefore acts not as an unbiased sample of the field, but as part of a dynamic reconstruction process governed by an internal geometric flow.

To describe this reconstruction mathematically, consider the free-energy functional
\[
\mathcal{F}[x] = \int \left( \frac{1}{2} |\dot{x}(t)|^{2} + S(t) \right) dt,
\]
where \(S(t)\) is the entropy field induced by quantization. Minimizing \(\mathcal{F}\) yields a field configuration that balances smoothness in \(x(t)\) with the entropy injected by measurement. This balance is analogous to Tikhonov regularization in inverse problems, in which a smoothness constraint is combined with data fidelity. The delta–sigma loop implements this minimization recurrently. The state variable \(x[n]\) evolves according to a discrete gradient descent:
\[
x[n+1] - x[n] = -\eta \frac{\delta \mathcal{F}}{\delta x[n]},
\]
where the derivative involves both the smoothness term and the entropy term. This evolution ensures that the resulting state trajectory closely matches the true field, even when the field is not bandlimited and even when the samples are corrupted by highly nonlinear collapse.

This procedure suggests a general theory of discrete field sampling that does not rely on bandlimitation. Instead, reconstruction is performed by a dynamical system whose vector flow approximates the curvature of the underlying field. Let the underlying field be a solution of some unknown differential equation, such as
\[
\frac{d^{L} \Phi}{dt^{L}} = f(\Phi, t),
\]
where \(f\) may encode physical laws or arbitrary dynamics. A delta–sigma modulator can be interpreted as an approximate solver for this equation. It estimates the curvature of \(\Phi\) by integrating discrepancies over time, and it corrects itself whenever the quantizer collapse introduces error. Thus the converter is capable of reconstructing \(\Phi\) even if \(\Phi\) is not bandlimited. What guarantees accuracy is not a frequency-domain constraint, but the stability of the underlying RSVP flow.

To express this reconstruction explicitly, one defines the effective operator
\[
\mathcal{T}_{\Delta}[\Phi] = \Phi - Q(\Phi),
\]
which measures the deviation between the field and its quantized representation. The delta–sigma loop applies the operator
\[
\mathcal{R}[\Phi] = \left(\frac{d}{dt}\right)^{-L} \mathcal{T}_{\Delta}[\Phi],
\]
which corresponds to integrating the discrepancy repeatedly. This operator acts as a smoothing operator or Green’s function associated with a high-order differential operator. Under appropriate stability conditions, \(\mathcal{R}\) converges to a projection onto a subspace of low-curvature fields. Thus the converter reconstructs the low-curvature component of \(\Phi\), effectively filtering out high-frequency components of the entropy field while preserving the slow variation of the true underlying field.

The sampling process therefore becomes a form of forced relaxation, in which the quantizer introduces perturbations that are subsequently dissipated through high-order integration. This is analogous to the lamphrodynamic processes described earlier, in which a field returns to a stable configuration after being disturbed. In delta–sigma modulation, the stable configuration is the estimate of the underlying field, while the disturbances are the entropy injections from quantization. The modulator repeatedly injects entropy, but the vector field routes it into high-frequency modes, where it is harmless.

This general theory of field sampling reveals that delta–sigma modulation can be interpreted as a computational process that maps a continuous field into a discrete representation while preserving structure. It is a special case of a far broader class of sampling systems that includes biological sensors, cognitive perceptual loops, and physical measurement systems. All such systems interact with a continuous field through discrete collapse operations and subsequently correct for the collapse through dynamical inference.

By interpreting sampling as a combination of collapse and inference, governed by the geometry of the error and entropy fields, we obtain a new set of design principles for creating converters that operate in nontraditional domains. For example, the theory applies to fields defined on manifolds, such as the curvature of a surface or the flow of a vector field in a mechanical system. It also applies to fields that evolve under stochastic dynamics, such as temperature fluctuations in metrological measurements. In all cases, delta–sigma-like systems can be designed to reconstruct these fields through discrete inference with entropy routing.

In summary, delta–sigma modulation is not merely an engineering technique for shaping quantization noise. It is a representative of a far more general theory of discrete field sampling in which measurement, collapse, inference, and curvature interact to produce stable reconstruction of continuous fields. This viewpoint provides the foundation for the next chapter, which develops future research directions and theoretical extensions of delta–sigma modulators as field-theoretic, geometric, and cognitive inference machines.


